% GENERATED FILE. Edit canonical lessons, then run npm run generate:books.
% canonical-source-hash: fnv1a64:55e15979971e5256
% canonical-lessons: MW-C157-koni, MW-C157-mat, MW-C157-samajhno, MW-W47-jha, MW-R157-not-today

\chapter{Not Today}
\label{ch:mw-a2-not-today}
\hlchaptermodality{157}

\begin{chapteropening}
\textbf{By the end of this chapter:} \emph{I can say that I or someone else does not do something, that someone is not here, and ask someone not to do something.}

Complete MW-R157-not-today: say what you and others do not do, who is not here, and ask someone not to go.
\end{chapteropening}

\section[\texorpdfstring{\mw{कोनी} (konī)}{कोनी (konī)}]{\mw{कोनी} — not}

\label{lesson:MW-C157-koni}



\begin{quote}
Say \emph{no}, the way you turn down a price: \textbf{\mw{कोनी।}} Then say \emph{I drink tea}: \textbf{\mw{हूं} \mw{चाय} \mw{पीऊं} \mw{हूं।}}
\end{quote}

\begin{grammarlens}[title={Grammar Lens: \mw{कोनी} before the verb}]
The word you refuse with also turns a sentence round. Put \textbf{\mw{कोनी}} in front of the verb, and let the \textbf{\mw{हूं}} or \textbf{\mw{है}} at the end go:

\begin{quote}
\textbf{\mw{हूं} \mw{चाय} \mw{कोनी} \mw{पीऊं।}}
\end{quote}

\begin{quote}
\emph{hū̃ chāy konī pīū̃.} — I don't drink tea.
\end{quote}

\begin{quote}
\textbf{\mw{वो} \mw{मांस} \mw{कोनी} \mw{खावै।}}
\end{quote}

\begin{quote}
\emph{vo mā̃s konī khāvai.} — He doesn't eat meat.
\end{quote}
\end{grammarlens}

\begin{grammarlens}[title={Grammar Lens: \mw{कोनी} in place of \mw{है}}]
With no other verb, \textbf{\mw{कोनी}} stands where \textbf{\mw{है}} stood:

\begin{quote}
\textbf{\mw{राम} \mw{आज} \mw{अठै} \mw{कोनी।}}
\end{quote}

\begin{quote}
\emph{rām āj aṭhai konī.} — Ram isn't here today.
\end{quote}
\end{grammarlens}

\subsection*{Your turn}
Say these aloud:

\begin{itemize}
\raggedright
  \item \textbf{\mw{हूं} \mw{मांस} \mw{कोनी} \mw{खाऊं।}} — \emph{hū̃ mā̃s konī khāū̃.} — I don't eat meat.
  \item \textbf{\mw{आप} \mw{चाय} \mw{कोनी} \mw{पीवो।}} — \emph{āp chāy konī pīvo.} — You don't drink tea.
  \item \textbf{\mw{दूध} \mw{गरम} \mw{कोनी।}} — \emph{dūdh garam konī.} — The milk isn't hot.
\end{itemize}

\begin{tcolorbox}[breakable,colback=teal!4,colframe=teal!35!black,title={Before you move on}]
Where does \textbf{\mw{कोनी}} go? (In front of the verb.) What goes from the end? (\textbf{\mw{हूं}} or \textbf{\mw{है}}.)
\end{tcolorbox}
\section[\texorpdfstring{\mw{मत} (mat)}{मत (mat)}]{\mw{मत} — don't}

\label{lesson:MW-C157-mat}



\begin{quote}
Ask for tea: \textbf{\mw{चाय} \mw{लावो।}} Then say \emph{I don't drink tea}.
\end{quote}

\begin{grammarlens}[title={Grammar Lens: \mw{मत} before a request}]
\textbf{\mw{मत}} in front of a request makes it \emph{don't}:

\begin{quote}
\textbf{\mw{मत} \mw{जावो।}}
\end{quote}

\begin{quote}
\emph{mat jāvo.} — Don't go.
\end{quote}

\begin{quote}
\textbf{\mw{चाय} \mw{मत} \mw{लावो।}}
\end{quote}

\begin{quote}
\emph{chāy mat lāvo.} — Don't bring tea.
\end{quote}

\textbf{\mw{कोनी}} says what is not so; \textbf{\mw{मत}} asks someone not to. Old Marwari had \emph{mati} in the same use (CDIAL 9981).
\end{grammarlens}

\subsection*{Your turn}
Say these aloud:

\begin{itemize}
\raggedright
  \item \textbf{\mw{अबार} \mw{मत} \mw{जावो।}} — \emph{abār mat jāvo.} — Don't go now.
  \item \textbf{\mw{मत} \mw{भूलो।}} — \emph{mat bhūlo.} — Don't forget.
  \item \textbf{\mw{हूं} \mw{चाय} \mw{कोनी} \mw{पीऊं।} \mw{चाय} \mw{मत} \mw{लावो।}} — \emph{hū̃ chāy konī pīū̃. chāy mat lāvo.} — I don't drink tea. Don't bring tea.
\end{itemize}

\begin{tcolorbox}[breakable,colback=teal!4,colframe=teal!35!black,title={Before you move on}]
\emph{Don't go}? (\textbf{\mw{मत} \mw{जावो।}}) \emph{I don't go}? (\textbf{\mw{हूं} \mw{कोनी} \mw{जाऊं।}})
\end{tcolorbox}
\section[\texorpdfstring{\mw{समझणो} (samajhṇo)}{समझणो (samajhṇo)}]{\mw{समझणो} (samajhṇo) — to understand}

\label{lesson:MW-C157-samajhno}



\begin{quote}
Before the new one: say \emph{to know}, \emph{to learn} and \emph{to think}. (\textbf{\mw{जाणणो}}, \textbf{\mw{सीखणो}}, \textbf{\mw{सोचणो}}.)
\end{quote}

\subsection*{You'll want to know: \mw{समझणो}}
\textbf{\mw{समझणो}} — \emph{samajhṇo} — \textquotedblleft{}to understand\textquotedblright{}.

\begin{quote}
\textbf{\mw{हूं} \mw{थोड़ु} \mw{समझूं} \mw{हूं।}}
\end{quote}

\begin{quote}
\emph{hū̃ thoṛu samajhū̃ hū̃.} — I understand a little.
\end{quote}

\begin{quote}
\textbf{\mw{हूं} \mw{कोनी} \mw{समझूं।}}
\end{quote}

\begin{quote}
\emph{hū̃ konī samajhū̃.} — I don't understand.
\end{quote}

It holds a letter you have not written yet: \emph{jh}, a \textbf{\mw{ज}} with a breath after it. The next lesson puts it in your hand.

\begin{cousinweb}
\textbf{\mw{समझणो}} (\emph{samajhṇo}) is inherited from Sanskrit \emph{sambudhyate}, \textquotedblleft{}wakes up, understands\textquotedblright{}. CDIAL lists Old Marwari \emph{samajhaï}, \textquotedblleft{}understands\textquotedblright{} (CDIAL 12959).
\end{cousinweb}

\begin{grammarlens}[title={What you've built}]
Its stem \textbf{\mw{समझ}} ends in a consonant, so it goes like \textbf{\mw{बोलणो}}: \textbf{\mw{समझूं} \mw{हूं}}, \textbf{\mw{समझै} \mw{है}}, \textbf{\mw{समझो} \mw{हो}}.
\end{grammarlens}

\subsection*{Your turn}
Say these aloud:

\begin{itemize}
\raggedright
  \item \emph{samajhṇo}
  \item \textbf{\mw{हूं} \mw{कोनी} \mw{समझूं।}} — \emph{hū̃ konī samajhū̃.} — I don't understand.
  \item \textbf{\mw{वो} \mw{थोड़ु} \mw{समझै} \mw{है।}} — \emph{vo thoṛu samajhai hai.} — He understands a little.
\end{itemize}

\begin{tcolorbox}[breakable,colback=teal!4,colframe=teal!35!black,title={Before you move on}]
What does \textbf{\mw{समझणो}} mean? (\textquotedblleft{}To understand\textquotedblright{}.) \emph{I don't understand}? (\textbf{\mw{हूं} \mw{कोनी} \mw{समझूं।}})
\end{tcolorbox}
\section[\texorpdfstring{\mw{झ} (jha)}{झ (jha)}]{\mw{झ} (jha) — the letter inside \mw{समझणो}}

\label{lesson:MW-W47-jha}



\begin{quote}
Before the new one: write \textbf{\mw{ज}} once, and name it.

Now say \textbf{\mw{समझणो}} (\emph{samajhṇo}), \textquotedblleft{}to understand\textquotedblright{}. Its middle is a \textbf{\mw{ज}} with a breath after it: \emph{jh}. This lesson puts that letter in your hand.
\end{quote}

\subsection*{Script you'll notice: \mw{झ}}
\textbf{\mw{झ}} — \emph{jha}, the breathed \emph{j}.

It sits inside \textbf{\mw{समझणो}}. It pairs with \textbf{\mw{ज}} as \textbf{\mw{छ}} pairs with \textbf{\mw{च}}: the same place in the mouth, with a breath. On the page it is one joined body, with a short upper stem, two loops and a diagonal tail, then a middle crossbar, a right stem and the headline.

\subsection*{Writing: \mw{झ}}
\hlblockfigure{\detokenize{figures/MW-W47-jha-filmstrip.pdf}}{How \mw{झ} is written, stroke by stroke}

\begin{itemize}
\raggedright
  \item \textbf{1.} descend the short upper stem, turn left and curve clockwise around the upper bowl, continue through the waist and clockwise around the lower loop, then finish down-right through the diagonal tail without lifting
  \item \textbf{2.} lift at the central junction and draw the middle crossbar left-to-right
  \item \textbf{3.} without lifting, climb up the right stem to the headline
  \item \textbf{4.} descend the right stem top-to-bottom
  \item \textbf{5.} lift and draw the top shirorekhā left-to-right
\end{itemize}

\textbf{Pen lifts: 2.}

\begin{quote}
The source shows four movements; native writers draw \mw{झ} as three pen-down strokes (53\% of HP Labs India's native-writer samples, 44 of 83). This ductus follows the four movements in order and direction, lifting only before the crossbar and the headline: the crossbar ends on the right stem, so the pen climbs the stem before descending it.
\end{quote}

\begin{quote}
Stroke order is one attested teaching order; hands and teachers vary. Source: Opiaterein, ‘Deva-\mw{झ}-order.gif’, strokes 1–4, Wikimedia Commons, 10 May 2009.
\end{quote}

\subsection*{Your turn}
\begin{itemize}
\raggedright
  \item \emph{Look:} at \textbf{\mw{समझणो}} and point at \mw{झ} inside it
  \item \emph{Trace it:} \mw{झ} three times, saying \emph{jh} as you finish each one
  \item \emph{Write it:} \mw{ज}, then \mw{झ}, from memory
  \item \emph{Read it:} \textbf{\mw{समझणो}} aloud once more
\end{itemize}

\begin{tcolorbox}[breakable,colback=teal!4,colframe=teal!35!black,title={Before you move on}]
Which letter is this — \mw{झ}? (\textbf{jha}, the breathed \emph{j}.) Which word did it come from? (\textbf{\mw{समझणो}}, \emph{samajhṇo}, \textquotedblleft{}to understand\textquotedblright{}.)
\end{tcolorbox}
\section[(dialogue)]{Not today — a review}

\label{lesson:MW-R157-not-today}



\begin{quote}
Say \emph{I don't drink tea}.
\end{quote}

\subsection*{Your turn}
\begin{itemize}
\raggedright
  \item \emph{Say it:} \emph{I don't eat meat}; \emph{he doesn't eat meat}; \emph{you don't drink tea}, to \textbf{\mw{आप}}
  \item \emph{Say it:} \emph{Ram isn't here today}; \emph{my cow isn't here}
  \item \emph{Say it:} \emph{don't go}; \emph{don't go now}; \emph{don't bring tea}
  \item \emph{Say it:} \emph{I understand a little}; \emph{I don't understand}
  \item \emph{Read it:} \textbf{\mw{समझणो}}, and name the letter after \textbf{\mw{म}}
\end{itemize}

\begin{tcolorbox}[breakable,colback=teal!4,colframe=teal!35!black,title={Before you move on}]
\emph{I don't understand}? (\textbf{\mw{हूं} \mw{कोनी} \mw{समझूं।}}) \emph{Ram isn't here}? (\textbf{\mw{राम} \mw{अठै} \mw{कोनी।}}) \emph{Don't go}? (\textbf{\mw{मत} \mw{जावो।}})
\end{tcolorbox}
